% =============================================================================
% OXIMETRIA DE PULSO — ONE PAGER ICU
% Versão funcional: cartão de referência rápida para beira-leito
% Estilo: denso, color-coded, prático (inspirado em cartões de emergência)
% Nick Mark MD — Tradução
% =============================================================================
\documentclass[landscape,a3paper]{article}
\usepackage[landscape,a3paper,margin=6mm]{geometry}
\usepackage[utf8]{inputenc}
\usepackage[T1]{fontenc}
\usepackage[brazilian]{babel}
\usepackage{lmodern}
\renewcommand{\familydefault}{\sfdefault}
\usepackage{microtype}
\usepackage{xcolor}
\usepackage{colortbl}
\usepackage{booktabs}
\usepackage{multicol}
\usepackage{array}
\usepackage{calc}
\usepackage{tikz}
\usepackage{amsmath}
\usetikzlibrary{arrows.meta}

\pagestyle{empty}
\setlength{\parindent}{0pt}
\setlength{\parskip}{0pt}

% ---------------------------------------------------------------------------
% CORES
% ---------------------------------------------------------------------------
\definecolor{titleBar}{RGB}{20, 40, 70}
\definecolor{physRed}{RGB}{165, 30, 35}
\definecolor{physLight}{RGB}{252, 228, 230}
\definecolor{sensorBlue}{RGB}{25, 80, 150}
\definecolor{sensorLight}{RGB}{215, 230, 248}
\definecolor{precGreen}{RGB}{30, 110, 55}
\definecolor{precLight}{RGB}{218, 242, 225}
\definecolor{hemPurple}{RGB}{90, 40, 130}
\definecolor{hemLight}{RGB}{235, 220, 248}
\definecolor{warnOrange}{RGB}{200, 120, 0}
\definecolor{warnYellow}{RGB}{255, 243, 200}
\definecolor{waveBlue}{RGB}{10, 90, 170}
\definecolor{waveLight}{RGB}{220, 235, 250}
\definecolor{stepGray}{RGB}{90, 90, 90}
\definecolor{zebraA}{RGB}{255, 255, 255}
\definecolor{zebraB}{RGB}{240, 245, 252}
\definecolor{oxyBlue}{RGB}{2, 132, 199}
\definecolor{deoxyRed}{RGB}{180, 30, 30}
\definecolor{irOrange}{RGB}{210, 140, 20}

% ---------------------------------------------------------------------------
% COMANDOS
% ---------------------------------------------------------------------------

% Cabeçalho de seção (barra colorida)
\newcommand{\secbar}[2]{%
  \par\vspace{3pt}%
  \noindent\colorbox{#1}{\parbox{\dimexpr\linewidth-4pt}{%
    \centering\footnotesize\bfseries\strut\textcolor{white}{#2}\strut%
  }}%
  \par\vspace{2pt}%
}

% Caixa de conteúdo
\newcommand{\infocard}[4]{%
  % #1=cor acento, #2=título, #3=conteúdo, #4=alerta (pode ser vazio)
  \par\vspace{2pt}%
  \noindent\fcolorbox{#1!70!black}{#1!6!white}{\parbox{\dimexpr\linewidth-6pt}{%
    \noindent\colorbox{#1}{\parbox{\dimexpr\linewidth-2pt}{%
      \scriptsize\bfseries\color{white}\strut\ #2\strut%
    }}\\[2pt]%
    \scriptsize%
    #3%
    \ifx&#4&\else\\[2pt]\noindent\colorbox{warnYellow}{\parbox{\dimexpr\linewidth-6pt}{\scriptsize\textcolor{warnOrange}{\bfseries !}\ #4}}\fi%
  }}%
  \par%
}

% Ponto-chave com bullet
\newcommand{\kp}[1]{{\color{stepGray}\bfseries\guillemotright}\hspace{0.4em}#1\par\vspace{1.5pt}}

% Valor de SpO2 em destaque
\newcommand{\spo}[1]{\textbf{\textcolor{oxyBlue}{#1}}}

\begin{document}

% === TITLE BAR ===
\noindent\colorbox{titleBar}{\parbox{\dimexpr\textwidth-4pt}{%
  \centering\vspace{2pt}%
  {\color{white}\bfseries\Large OXIMETRIA DE PULSO}%
  \quad{\color{white}\footnotesize\textbar\quad Princípios, Precisão e Análise da Onda Pletismográfica}%
  \hfill{\color{white!70}\scriptsize One Pager ICU $\bullet$ Nick Mark MD $\bullet$ Tradução}%
  \vspace{2pt}%
}}

\vspace{2pt}

% === LEGENDA ===
\noindent\colorbox{sensorLight}{\parbox{\dimexpr\textwidth-4pt}{%
  \centering\scriptsize%
  \textbf{Como ler este cartão:}
  \textcolor{physRed}{$\blacksquare$}\,Princípio
  \quad\textcolor{sensorBlue}{$\blacksquare$}\,Sensor/Técnica
  \quad\textcolor{precGreen}{$\blacksquare$}\,Fatores de Precisão
  \quad\textcolor{hemPurple}{$\blacksquare$}\,Hemoglobinas\,\&\,Interferências
  \quad\textcolor{waveBlue}{$\blacksquare$}\,Onda Pletismográfica
  \quad\textcolor{warnOrange}{$\blacksquare$}\,Alertas
}}

\vspace{3pt}

\setlength{\columnsep}{6pt}
\setlength{\columnseprule}{0.3pt}
\renewcommand{\columnseprulecolor}{\color{gray!40}}

\begin{multicols}{4}
\scriptsize
\setlength{\tabcolsep}{2pt}
\renewcommand{\arraystretch}{1.1}

% =====================================================================
% COLUNA 1 — PRINCÍPIO E SENSOR
% =====================================================================

\secbar{physRed}{PRINCÍPIO FISIOLÓGICO}

\infocard{physRed}{LEI DE BEER-LAMBERT}{%
    A oximetria de pulso mede a \textbf{saturação de O\textsubscript{2} da hemoglobina} de forma contínua e não invasiva.\\[2pt]
    Explora a \textbf{absorção diferencial} de luz entre:\\[1pt]
    \kp{\textcolor{oxyBlue}{\textbf{OxyHb}} --- absorve mais \textcolor{irOrange}{IR (940\,nm)}}
    \kp{\textcolor{deoxyRed}{\textbf{DeoxyHb}} --- absorve mais \textcolor{deoxyRed}{Red (660\,nm)}}
    A razão \textbf{R = (AC/DC)\textsubscript{Red} / (AC/DC)\textsubscript{IR}} é convertida em SpO\textsubscript{2} usando uma \textbf{curva de calibração empírica}.
}{}

\infocard{physRed}{COMPONENTES DO SINAL}{%
    \begin{tabular}{@{}>{\bfseries}l p{4.5cm}@{}}
    \rowcolor{physLight}
    AC (pulsátil) & Sangue \textbf{arterial} apenas. Varia com o ciclo cardíaco. \textbf{É isso que o oxímetro usa.} \\
    DC (constante) & Tecidos, ossos, sangue venoso e capilar. Filtrado pelo algoritmo. \\
    \end{tabular}
}{}

\vspace{2pt}

% GRÁFICO: Curvas de Absorção
\noindent\fcolorbox{physRed!70!black}{physRed!6!white}{\parbox{\dimexpr\linewidth-6pt}{%
  \noindent\colorbox{physRed}{\parbox{\dimexpr\linewidth-2pt}{%
    \scriptsize\bfseries\color{white}\strut\ ESPECTRO DE ABSORÇÃO\strut%
  }}\\[3pt]%
  \centering
  \begin{tikzpicture}[scale=0.62, font=\sffamily\tiny]
      \draw[->, stepGray] (0,0) -- (7,0) node[right] {$\lambda$ (nm)};
      \draw[->, stepGray] (0,0) -- (0,4.5) node[above] {Absorção};

      % Ticks
      \foreach \x/\t in {0.5/600, 2/700, 3.5/800, 5/900, 6.5/1000} {
          \draw[stepGray] (\x, 0) -- (\x, -0.1) node[below, scale=0.9] {\t};
      }

      % Curvas
      \draw[ultra thick, oxyBlue] (0.3, 4) .. controls (1.5, 3.5) and (2.5, 2.2) ..
          (3.5, 2.2) .. controls (4.5, 2.2) and (5.5, 1.5) .. (6.5, 1.2);
      \node[oxyBlue, right, scale=0.9] at (6.5, 1.2) {\textbf{OxyHb}};

      \draw[ultra thick, deoxyRed] (0.3, 1) .. controls (1, 2) and (2, 3.5) ..
          (3.5, 3.5) .. controls (5, 3.5) and (5.8, 2.5) .. (6.5, 2.5);
      \node[deoxyRed, right, scale=0.9] at (6.5, 2.5) {\textbf{DeoxyHb}};

      % Linhas de referência
      \draw[dashed, deoxyRed!60, thick] (1.0, 0) -- (1.0, 4.3);
      \node[above, deoxyRed, scale=0.85, font=\bfseries] at (1.0, 4.3) {660};

      \draw[dashed, irOrange, thick] (5.5, 0) -- (5.5, 4.3);
      \node[above, irOrange, scale=0.85, font=\bfseries] at (5.5, 4.3) {940};
  \end{tikzpicture}
  \par
}}
\par

\vspace{2pt}

\secbar{sensorBlue}{SENSOR E POSICIONAMENTO}

\infocard{sensorBlue}{TIPOS DE SENSOR}{%
    \begin{tabular}{@{}>{\bfseries}l p{4.0cm}@{}}
    \rowcolor{sensorLight}
    Transmissão & Dedo, pé (neonato). LED de um lado, detector do outro. \textbf{Mais preciso.} \\
    Reflexão & Testa, orelha. LED e detector do mesmo lado. Útil quando não há acesso aos dedos. \\
    \end{tabular}\\[2pt]
    \kp{\textbf{Dedo} --- mais preciso. Polegar pode ser marginalmente superior.}
    \kp{\textbf{Orelha\,\&\,Testa} --- mais rápidos e confiáveis em \textbf{vasoconstrição} e hipotermia.}
    \kp{Nenhum local é \textit{sempre} superior --- teste vários.}
}{Posicionar com LED sobre a unha. Evitar edema, curativos e luz ambiente direta.}

% --- DIAGRAMA DO SENSOR (compacto) ---
\noindent\fcolorbox{sensorBlue!70!black}{sensorBlue!6!white}{\parbox{\dimexpr\linewidth-6pt}{%
  \noindent\colorbox{sensorBlue}{\parbox{\dimexpr\linewidth-2pt}{%
    \scriptsize\bfseries\color{white}\strut\ SENSOR DE TRANSMISSÃO\strut%
  }}\\[3pt]%
  \centering
  \begin{tikzpicture}[scale=0.9, font=\sffamily\tiny]
      % Dedo
      \draw[thick, stepGray!50, fill=sensorLight!30] (-2, -0.6) .. controls (1, -0.6) and (1.4, -0.4) .. (1.6, 0) .. controls (1.4, 0.4) and (1, 0.6) .. (-2, 0.6);
      \draw[stepGray!30] (0.7, 0.4) .. controls (1.4, 0.4) and (1.6, 0.25) .. (1.6, 0) .. controls (1.6, -0.25) and (1.4, -0.4) .. (0.7, -0.4);

      % Sensor clamp
      \fill[sensorBlue!80] (-0.6, 0.9) rectangle (0.8, 1.15);
      \fill[sensorBlue!80] (-0.6, -0.9) rectangle (0.8, -1.15);
      \node[white, font=\bfseries\tiny] at (0.1, 1.02) {LED};
      \node[white, font=\bfseries\tiny] at (0.1, -1.02) {FOTO};

      % Raios
      \draw[dashed, deoxyRed, -{Latex[length=2mm]}, thick] (-0.05, 0.85) -- (-0.05, -0.85);
      \draw[dashed, irOrange, -{Latex[length=2mm]}, thick] (0.3, 0.85) -- (0.3, -0.85);
      \node[right, deoxyRed, scale=0.85] at (-0.05, 0.4) {R};
      \node[right, irOrange, scale=0.85] at (0.3, -0.4) {IR};
  \end{tikzpicture}
  \par\vspace{1pt}
}}
\par

% =====================================================================
% COLUNA 2 — FATORES DE PRECISÃO
% =====================================================================
\columnbreak

\secbar{precGreen}{FATORES QUE AFETAM A PRECISÃO}

\infocard{precGreen}{COR DA PELE}{%
    A melanina absorve luz e pode causar \textbf{superestimação} da SpO\textsubscript{2} real.\\[2pt]
    \begin{tabular}{@{}p{\dimexpr\linewidth-4pt}@{}}
    \rowcolor{precLight}
    Pacientes negros: \textbf{3$\times$ mais risco de hipoxemia oculta} (SpO\textsubscript{2} aparentemente normal com SaO\textsubscript{2} real baixa). \\
    \end{tabular}\\[2pt]
    \kp{Diferença pode chegar a \textbf{4--8\%} em hipoxemia moderada.}
    \kp{Sensor de reflexão (testa) pode atenuar parcialmente.}
}{A oximetria pode ser \textbf{falsamente tranquilizadora} em pele negra. Dosar SaO\textsubscript{2} na gasometria se suspeitar.}

\infocard{precGreen}{ESMALTE DE UNHA}{%
    Esmaltes absorvem luz adicional no espectro vermelho/IR.\\[2pt]
    \begin{tabular}{@{}>{\bfseries\scriptsize}l >{\scriptsize}l@{}}
    \rowcolor{precLight}
    Cores que interferem & Azul, verde, preto \\
    Cores seguras & Vermelho, rosa, francesinha \\
    \rowcolor{precLight}
    Conduta & Remover ou usar outro dedo \\
    \end{tabular}
}{}

\infocard{precGreen}{ESTADOS DE BAIXO FLUXO}{%
    \kp{\textbf{Choque} --- vasoconstrição reduz componente AC.}
    \kp{\textbf{Hipotensão} --- sinal fraco e leitura instável.}
    \kp{\textbf{ECMO\,/\,LVAD} --- fluxo não-pulsátil. Oximetria convencional \textbf{inútil}.}
    \kp{\textbf{Hipotermia} --- vasoconstrição periférica. Usar probe central (testa).}
}{Índice de Perfusão (PI) $<$ 0,4\% sugere sinal inadequado. Se PI baixo, desconfiar da leitura.}

\infocard{precGreen}{PRECISÃO EM BAIXA SpO\textsubscript{2}}{%
    A curva de calibração foi derivada de \textbf{voluntários saudáveis} (dados confiáveis até $\sim$75\%).\\[2pt]
    \begin{tabular}{@{}>{\scriptsize}l >{\scriptsize}l@{}}
    \rowcolor{precLight}
    \textbf{SpO\textsubscript{2} $>$ 80\%} & Precisão $\pm$ 2\% \\
    \textbf{SpO\textsubscript{2} 70--80\%} & Precisão $\pm$ 3--5\% \\
    \rowcolor{precLight}
    \textbf{SpO\textsubscript{2} $<$ 70\%} & \textbf{Não confiável} \\
    \end{tabular}
}{}

\infocard{precGreen}{HIPEROXEMIA}{%
    A oximetria não distingue PaO\textsubscript{2} normal de supra-normal quando SpO\textsubscript{2} bate em 100\%.\\[2pt]
    \kp{Alvo razoável: \spo{SpO\textsubscript{2} 94--98\%}}
    \kp{Pós-parada cardíaca: evitar PaO\textsubscript{2} $>$ 300\,mmHg.}
}{Hiperóxia causa vasoconstrição cerebral, lesão oxidativa e piora desfechos pós-PCR.}

\infocard{precGreen}{ATRASO NA LEITURA (LAG)}{%
    \begin{tabular}{@{}>{\bfseries\scriptsize}l >{\scriptsize}l@{}}
    \rowcolor{precLight}
    Dedo & 5--15\,s de atraso \\
    Orelha/Testa & 2--5\,s (mais central) \\
    \rowcolor{precLight}
    Hipotensão & Lag prolongado \\
    \end{tabular}\\[2pt]
    \kp{A SpO\textsubscript{2} pode \textbf{continuar caindo} após intubação correta.}
    \kp{Não usar como única confirmação de via aérea.}
}{}

% =====================================================================
% COLUNA 3 — HEMOGLOBINAS + INTERFERÊNCIAS
% =====================================================================
\columnbreak

\secbar{hemPurple}{HEMOGLOBINAS ANORMAIS}

\infocard{hemPurple}{CARBOXIHEMOGLOBINA (CoHb)}{%
    O oxímetro convencional \textbf{não distingue} CoHb de OxyHb (ambas absorvem de forma semelhante no Red).\\[2pt]
    \begin{tabular}{@{}>{\bfseries\scriptsize}l >{\scriptsize}l@{}}
    \rowcolor{hemLight}
    Leitura & \textbf{Falsamente NORMAL} (90--100\%) \\
    Realidade & Paciente profundamente hipóxico \\
    \rowcolor{hemLight}
    Quando & Incêndio, escapamento, narguile \\
    Diagnóstico & Co-oximetria (SaO\textsubscript{2} real) \\
    \end{tabular}
}{SpO\textsubscript{2} ``normal'' em vítima de incêndio NÃO exclui intoxicação por CO!}

\infocard{hemPurple}{METEMOGLOBINEMIA (MetHb)}{%
    MetHb absorve igualmente Red e IR $\rightarrow$ razão R $\rightarrow$ 1 $\rightarrow$ SpO\textsubscript{2} trava em $\sim$\textbf{85\%}.\\[2pt]
    \begin{tabular}{@{}>{\bfseries\scriptsize}l >{\scriptsize}l@{}}
    \rowcolor{hemLight}
    Leitura & ``Trava'' em $\sim$85\% \\
    Causas & Dapsona, anestésicos locais, nitratos \\
    \rowcolor{hemLight}
    Clínica & Cianose com PaO\textsubscript{2} normal \\
    Conduta & Co-oximetria + azul de metileno \\
    \end{tabular}
}{MetHb $>$ 1\%: leitura não confiável. SpO\textsubscript{2} converge para 85\% independente da SaO\textsubscript{2}.}

\infocard{hemPurple}{OUTRAS HEMOGLOBINAS}{%
    \kp{\textbf{Sulfa-Hb} --- leitura espuriamente \textbf{baixa}. Paciente pode não estar hipóxico.}
    \kp{\textbf{HbA1c $>$ 7} --- leve superestimativa da SpO\textsubscript{2}.}
    \kp{\textbf{Hemoglobina Fetal (HbF)} --- não afeta significativamente.}
}{Co-oximetria de pulso (4+ comprimentos de onda) mede MetHb e CoHb diretamente.}

\vspace{2pt}
\secbar{warnOrange}{INTERFERÊNCIAS TÉCNICAS}

\infocard{warnOrange}{CORANTES E MEDICAMENTOS IV}{%
    \begin{tabular}{@{}>{\bfseries\scriptsize}l >{\scriptsize}l >{\scriptsize}l@{}}
    \rowcolor{warnYellow}
    Azul de metileno & $\downarrow\downarrow$ SpO\textsubscript{2} (falsa) & 1--2\,min \\
    Indocianina verde & $\downarrow$ SpO\textsubscript{2} (falsa) & $<$ 5\,min \\
    \rowcolor{warnYellow}
    Isossulfan azul & $\downarrow$ SpO\textsubscript{2} (falsa) & minutos \\
    \end{tabular}\\[2pt]
    \kp{Efeito \textbf{transitório} --- esperar washout antes de confiar na leitura.}
}{}

\infocard{warnOrange}{ARTEFATO DE MOVIMENTO}{%
    \kp{Movimentos do paciente geram ondas falsas que mimetizam AC.}
    \kp{Tecnologias modernas (Masimo Signal Extraction, Nellcor OxiMax) filtram melhor.}
    \kp{Se SpO\textsubscript{2} oscila muito $\rightarrow$ provavelmente artefato.}
}{}

\infocard{warnOrange}{LUZ AMBIENTE}{%
    \kp{Luz fluorescente, cirúrgica ou solar direta pode interferir.}
    \kp{Cobrir o sensor com material opaco se necessário.}
}{}

% =====================================================================
% COLUNA 4 — ONDA PLETISMOGRÁFICA + SUMÁRIO
% =====================================================================
\columnbreak

\secbar{waveBlue}{ANÁLISE DA ONDA PLETISMOGRÁFICA}

\infocard{waveBlue}{ANATOMIA DA ONDA}{%
    \centering
    \begin{tikzpicture}[scale=0.75, font=\sffamily\tiny]
        \fill[waveLight] (0,0) .. controls (0.2,0) and (0.5,2.4) .. (0.8,2.4)
            .. controls (1.0,2.4) and (1.2,1) .. (1.5,1)
            .. controls (1.8,1.5) and (2.6,0) .. (3.2,0) -- cycle;
        \draw[thick, waveBlue] (0,0) .. controls (0.2,0) and (0.5,2.4) .. (0.8,2.4)
            .. controls (1.0,2.4) and (1.2,1) .. (1.5,1)
            .. controls (1.8,1.5) and (2.6,0) .. (3.2,0);

        \draw[dashed, stepGray!50] (0.8,0) -- (0.8,2.4);
        \draw[dashed, stepGray!50] (1.5,0) -- (1.5,1);

        \node[above, waveBlue, font=\bfseries\tiny] at (0.8, 2.4) {Pico Sistólico};
        \node[right, stepGray, font=\tiny] at (1.5, 1.1) {Nó Dicrótico};

        \draw[<->, stepGray] (0.8, -0.3) -- (1.5, -0.3) node[midway, below, font=\tiny] {$\Delta$T};
    \end{tikzpicture}
    \par\vspace{2pt}
    \flushleft
    \kp{\textbf{Pico Sistólico} = ejeção ventricular}
    \kp{\textbf{Nó Dicrótico} = fechamento da valva aórtica}
    \kp{\textbf{$\Delta$T} reflete compliance vascular}
}{}

\infocard{waveBlue}{PADRÕES CLÍNICOS}{%
    \colorbox{waveLight}{\parbox{\dimexpr\linewidth-4pt}{\centering\bfseries\scriptsize NORMAL --- onda bem definida}}\\[2pt]
    \centering
    \begin{tikzpicture}[scale=0.5]
        \draw[thick, waveBlue] (0,0) .. controls (0.2,0) and (0.4,1) .. (0.5,1) .. controls (0.6,1) and (0.7,0.4) .. (0.8,0.4) .. controls (0.9,0.8) and (1.5,0) .. (2,0)
             (2,0) .. controls (2.2,0) and (2.4,1) .. (2.5,1) .. controls (2.6,1) and (2.7,0.4) .. (2.8,0.4) .. controls (2.9,0.8) and (3.5,0) .. (4,0)
             (4,0) .. controls (4.2,0) and (4.4,1) .. (4.5,1) .. controls (4.6,1) and (4.7,0.4) .. (4.8,0.4) .. controls (4.9,0.8) and (5.5,0) .. (6,0);
    \end{tikzpicture}
    \par\vspace{3pt}

    \flushleft
    \colorbox{warnYellow}{\parbox{\dimexpr\linewidth-4pt}{\centering\bfseries\scriptsize SINAL FRACO --- $\downarrow$ perfusão}}\\[2pt]
    \centering
    \begin{tikzpicture}[scale=0.5]
        \draw[thick, stepGray] (0,0) .. controls (0.2,0) and (0.4,0.25) .. (0.5,0.25) .. controls (0.6,0.25) and (0.7,0.1) .. (0.8,0.1) .. controls (0.9,0.15) and (1.5,0) .. (2,0)
             (2,0) .. controls (2.2,0) and (2.4,0.25) .. (2.5,0.25) .. controls (2.6,0.25) and (2.7,0.1) .. (2.8,0.1) .. controls (2.9,0.15) and (3.5,0) .. (4,0)
             (4,0) .. controls (4.2,0) and (4.4,0.25) .. (4.5,0.25) .. controls (4.6,0.25) and (4.7,0.1) .. (4.8,0.1) .. controls (4.9,0.15) and (5.5,0) .. (6,0);
    \end{tikzpicture}
    \par\vspace{3pt}

    \flushleft
    \colorbox{physLight}{\parbox{\dimexpr\linewidth-4pt}{\centering\bfseries\scriptsize VASOCONSTRIÇÃO --- $\uparrow$ dicrótico, $\Delta$T curto}}\\[2pt]
    \centering
    \begin{tikzpicture}[scale=0.5]
        \draw[thick, physRed] (0,0) .. controls (0.2,0) and (0.4,1) .. (0.5,1) .. controls (0.6,1) and (0.7,0.7) .. (0.8,0.7) .. controls (0.9,0.9) and (1.5,0) .. (2,0)
             (2,0) .. controls (2.2,0) and (2.4,1) .. (2.5,1) .. controls (2.6,1) and (2.7,0.7) .. (2.8,0.7) .. controls (2.9,0.9) and (3.5,0) .. (4,0)
             (4,0) .. controls (4.2,0) and (4.4,1) .. (4.5,1) .. controls (4.6,1) and (4.7,0.7) .. (4.8,0.7) .. controls (4.9,0.9) and (5.5,0) .. (6,0);
    \end{tikzpicture}
    \par\vspace{3pt}

    \flushleft
    \colorbox{hemLight}{\parbox{\dimexpr\linewidth-4pt}{\centering\bfseries\scriptsize VASODILATAÇÃO --- $\uparrow$ sistólico, $\Delta$T longo}}\\[2pt]
    \centering
    \begin{tikzpicture}[scale=0.5]
        \draw[thick, hemPurple] (0,0) .. controls (0.2,0) and (0.4,1.6) .. (0.5,1.6) .. controls (0.6,1.6) and (0.7,0.2) .. (0.8,0.2) .. controls (0.9,0.4) and (1.5,0) .. (2,0)
             (2,0) .. controls (2.2,0) and (2.4,1.6) .. (2.5,1.6) .. controls (2.6,1.6) and (2.7,0.2) .. (2.8,0.2) .. controls (2.9,0.4) and (3.5,0) .. (4,0)
             (4,0) .. controls (4.2,0) and (4.4,1.6) .. (4.5,1.6) .. controls (4.6,1.6) and (4.7,0.2) .. (4.8,0.2) .. controls (4.9,0.4) and (5.5,0) .. (6,0);
    \end{tikzpicture}
    \par
}{}

\vspace{2pt}
\secbar{titleBar}{REGRA DE OURO}

\noindent\fcolorbox{titleBar}{warnYellow}{\parbox{\dimexpr\linewidth-6pt}{%
    \centering\scriptsize\bfseries
    Sempre correlacione:\\[1pt]
    \textbf{Onda} (forma) + \textbf{PI} (força) + \textbf{Pulso palpável}\\[2pt]
    Em dúvida: \textcolor{physRed}{\textbf{Gasometria Arterial (SaO\textsubscript{2})}}\\
    é o padrão-ouro mandatório.
}}

\vspace{3pt}

\noindent\fcolorbox{sensorBlue}{sensorLight}{\parbox{\dimexpr\linewidth-6pt}{%
    \scriptsize\bfseries\centering
    CHECKLIST RÁPIDO\\[1pt]
    \normalfont\scriptsize
    \begin{tabular}{@{}l p{4.5cm}@{}}
    \rowcolor{sensorLight}
    \textbf{SpO\textsubscript{2} baixa?} & Checar sensor, posição, PI \\
    \textbf{Onda ruim?} & Trocar dedo/local, aquecer \\
    \rowcolor{sensorLight}
    \textbf{Incêndio?} & Ignorar SpO\textsubscript{2}, pedir CoHb \\
    \textbf{Cianose + SpO\textsubscript{2} ok?} & Pensar MetHb \\
    \rowcolor{sensorLight}
    \textbf{Pós-IOT?} & Lembrar do lag (5--15s) \\
    \end{tabular}
}}

\end{multicols}

\vspace{2pt}

\noindent\colorbox{titleBar}{\parbox{\dimexpr\textwidth-4pt}{%
  \centering\vspace{1pt}%
  {\color{white!70}\scriptsize
  Oxímetros convencionais usam 2 comprimentos de onda $\rightarrow$ só detectam OxyHb e DeoxyHb.
  \quad$\bullet$\quad
  Co-oxímetros de pulso (ex:\ Masimo Rainbow) usam 7+ $\lambda$ $\rightarrow$ medem SpCO, SpMet, SpHb.
  \quad$\bullet$\quad
  Material exclusivamente educacional.}%
  \vspace{1pt}%
}}

\end{document}
